\documentclass[11pt,a4paper]{article}

\usepackage{amsmath,amssymb,amsthm,mathrsfs}
\usepackage{geometry}
\usepackage{hyperref}
\usepackage{enumitem}
\geometry{margin=1in}

\title{Schr\"odinger Evolution with Time-Dependent Self-Adjoint Extensions\\
and Associated Quantum Instruments}
\author{}
\date{}

\newtheorem{theorem}{Theorem}[section]
\newtheorem{proposition}[theorem]{Proposition}
\newtheorem{definition}[theorem]{Definition}
\newtheorem{assumption}[theorem]{Assumption}
\newtheorem{remark}[theorem]{Remark}

\begin{document}

\maketitle

\begin{abstract}
We present a rigorous operator-theoretic treatment of the Schr\"odinger equation
with time-dependent boundary conditions. Such systems are described as
time-dependent self-adjoint extensions of a fixed symmetric operator.
Using quadratic form methods we establish existence and uniqueness of the
unitary propagator. We then construct an associated positive operator-valued
measure (POVM) describing boundary energy transfer and prove that the
resulting structure defines a quantum instrument.
\end{abstract}

\section{Introduction}

Let $\Omega \subset \mathbb{R}^d$ be bounded with smooth boundary.
Consider the symmetric operator
\[
H_0 = -\frac{\hbar^2}{2m}\Delta
\quad \text{on} \quad C_c^\infty(\Omega)\subset L^2(\Omega).
\]
Self-adjoint realizations correspond to boundary conditions selecting
a self-adjoint extension in the sense of von Neumann.

We study families $\{H(t)\}_{t\in[0,T]}$ of self-adjoint extensions
with time-dependent boundary conditions and establish
their relation to quantum measurement theory.

\section{Quadratic Form Framework}

Let $\mathcal H = L^2(\Omega)$.

\begin{assumption}[Regularity]
\begin{enumerate}[label=(\roman*)]
\item For each $t$, $H(t)$ is self-adjoint.
\item There exists a common dense quadratic form domain $Q\subset\mathcal H$.
\item The associated closed quadratic forms $q_t$ satisfy
\[
t\mapsto q_t(\psi,\phi)
\]
is $C^1$ for all $\psi,\phi\in Q$.
\end{enumerate}
\end{assumption}

This includes time-dependent Robin boundary conditions and moving
boundaries after reduction to fixed domain.

\section{Existence of Unitary Propagator}

\begin{theorem}[Existence and uniqueness]
Under Assumption 2.1, there exists a unique strongly continuous
unitary propagator $U(t,s)$ satisfying
\begin{align*}
U(s,s)&=\mathbf 1,\\
U(t,r)U(r,s)&=U(t,s),\\
i\hbar \frac{d}{dt}U(t,s)\psi &= H(t)U(t,s)\psi
\quad (\psi\in Q).
\end{align*}
\end{theorem}

\begin{proof}
Each $q_t$ is closed and semibounded.
The map $t\mapsto q_t$ is $C^1$ on the common domain $Q$.
Hence Kisy\'nski's theorem on evolution equations with
time-dependent quadratic forms applies.
See Kato's monograph for detailed exposition.
\end{proof}

\section{Reduction to Fixed Domain}

Let $W(t)$ be a differentiable family of unitary operators such that
\[
W(t)D(H(t)) = D(H(0)).
\]
Define
\[
\widetilde H(t)
=
W(t)H(t)W(t)^{-1}
-
i\hbar \dot W(t)W(t)^{-1}.
\]

\begin{proposition}
The transformed state $\phi(t)=W(t)\psi(t)$ satisfies
\[
i\hbar \partial_t \phi(t)
=
\widetilde H(t)\phi(t)
\]
on fixed domain $D(H(0))$.
\end{proposition}

\begin{proof}
Direct differentiation and substitution of the Schr\"odinger equation.
\end{proof}

\section{Boundary Energy Transfer as POVM}

Define
\[
M(\Delta)
=
\int_\Delta
U(t,0)^\dagger
\partial_t H(t)
U(t,0)
\, dt.
\]

\begin{theorem}
If $\partial_t q_t$ is positive as quadratic form, then
$M$ defines a POVM on $[0,T]$.
\end{theorem}

\begin{proof}
For $\psi\in Q$,
\[
\langle \psi, M(\Delta)\psi\rangle
=
\int_\Delta
\partial_t q_t(U(t,0)\psi)\, dt.
\]
Positivity of $\partial_t q_t$ implies positivity of $M(\Delta)$.
Countable additivity follows from properties of the integral.
\end{proof}

\section{Associated Quantum Instrument}

\begin{definition}
A quantum instrument on $(X,\Sigma)$ is a $\sigma$-additive map
\[
\mathcal I:\Sigma\to \mathrm{CP}(\mathcal T(\mathcal H))
\]
such that $\mathcal I_X$ is trace-preserving.
\end{definition}

Define formally
\[
\mathcal I_\Delta(\rho)
=
\int_\Delta
K(t)\rho K(t)^\dagger dt,
\quad
K(t)=\sqrt{\partial_t H(t)}\, U(t,0).
\]

\begin{theorem}
$\mathcal I$ is a completely positive instrument whose
associated POVM equals $M$.
\end{theorem}

\begin{proof}
Each integrand is CP.
Trace identity:
\[
\mathrm{Tr}[\mathcal I_\Delta(\rho)]
=
\mathrm{Tr}[\rho M(\Delta)].
\]
\end{proof}

\section{Boundary--Measurement Equivalence}

\begin{theorem}
Let $H(t)$ be a $C^1$ family of self-adjoint extensions of a symmetric
operator $H_0$ satisfying Assumption 2.1.
Then:
\begin{enumerate}
\item Schr\"odinger evolution is unitary.
\item Domain variation induces a natural POVM $M$.
\item There exists an associated quantum instrument $\mathcal I$.
\item The geometric generator $-i\hbar \dot W W^{-1}$ equals
the Hamiltonian term of the minimal Stinespring dilation.
\end{enumerate}
\end{theorem}

\begin{proof}
Combine previous results and Stinespring's dilation theorem.
\end{proof}

\section{Worked Example: Particle in a Moving Interval}

We consider a quantum particle confined to a time-dependent interval
\[
\Omega(t) = [0,L(t)],
\quad L \in C^2([0,T]), \quad L(t) > 0.
\]

At each time $t$, define the Dirichlet Laplacian
\[
H(t) = -\frac{\hbar^2}{2m} \frac{d^2}{dx^2}
\]
with domain
\[
D(H(t))
=
\left\{
\psi \in H^2(0,L(t))
:
\psi(0)=\psi(L(t))=0
\right\}.
\]

\subsection{Self-Adjointness}

For each fixed $t$, $H(t)$ is self-adjoint on $L^2(0,L(t))$.
The associated quadratic form is
\[
q_t(\psi)
=
\frac{\hbar^2}{2m}
\int_0^{L(t)}
|\psi'(x)|^2 dx
\]
with form domain $H_0^1(0,L(t))$.

\subsection{Reduction to a Fixed Hilbert Space}

Define the unitary scaling operator
\[
(W(t)\psi)(y)
=
\sqrt{L(t)}\,\psi(L(t)y),
\quad y\in[0,1].
\]

Then
\[
W(t): L^2(0,L(t)) \to L^2(0,1)
\]
is unitary.

Define
\[
\phi(y,t)=W(t)\psi(x,t).
\]

\subsection{Transformed Generator}

A direct computation yields

\[
W(t) H(t) W(t)^{-1}
=
\frac{1}{L(t)^2}
\left(
-\frac{\hbar^2}{2m}\frac{d^2}{dy^2}
\right).
\]

Differentiating $W(t)$ gives

\[
\dot W(t) W(t)^{-1}
=
\frac{\dot L(t)}{2L(t)}
+
\frac{\dot L(t)}{L(t)} y \partial_y.
\]

Therefore the effective Hamiltonian on $L^2(0,1)$ is

\[
\widetilde H(t)
=
\frac{1}{L(t)^2}
\frac{p_y^2}{2m}
-
\frac{\dot L(t)}{2L(t)}
( y p_y + p_y y ),
\]

where
\[
p_y = -i\hbar \frac{d}{dy}.
\]

The operator
\[
D = \frac{1}{2}(y p_y + p_y y)
\]
is the self-adjoint generator of dilations on $H_0^1(0,1)$.

Thus

\[
\widetilde H(t)
=
\frac{1}{L(t)^2} H_D
-
\frac{\dot L(t)}{L(t)} D.
\]

\subsection{Energy Balance}

Let $\phi(t)$ solve
\[
i\hbar \partial_t \phi = \widetilde H(t)\phi.
\]

Then

\[
\frac{d}{dt}
\langle \phi, \widetilde H(t)\phi \rangle
=
\left\langle
\phi,
\partial_t \widetilde H(t)
\phi
\right\rangle.
\]

A direct computation gives

\[
\partial_t \widetilde H(t)
=
-2 \frac{\dot L}{L^3} H_D
-
\frac{\ddot L L - \dot L^2}{L^2} D.
\]

Thus energy variation arises entirely from boundary motion.

\subsection{Induced POVM}

Assume $\dot L(t)\ge 0$ for $t\in\Delta$.
Then the quadratic form derivative satisfies positivity:

\[
\partial_t q_t(\psi)
=
\frac{\hbar^2}{2m}
\dot L(t)
|\psi'(L(t))|^2.
\]

Hence define

\[
M(\Delta)
=
\int_\Delta
U(t,0)^\dagger
\partial_t H(t)
U(t,0)
dt.
\]

For $\psi$ in the instantaneous domain,

\[
\langle \psi, M(\Delta)\psi \rangle
=
\frac{\hbar^2}{2m}
\int_\Delta
\dot L(t)
|\partial_x \psi(L(t),t)|^2
dt.
\]

This explicitly shows:

\begin{itemize}
\item The POVM measures boundary gradient amplitude.
\item The statistics quantify probability/energy transfer
      induced by wall motion.
\item The dilation generator $D$ is precisely the
      backaction Hamiltonian.
\end{itemize}

\subsection{Adiabatic Regime}

If
\[
\left|\frac{\dot L}{L}\right|
\ll
\frac{E_n(t)-E_m(t)}{\hbar},
\]
transitions between instantaneous eigenstates
\[
\phi_n(y,t) = \sqrt{2}\sin(n\pi y)
\]
are suppressed.

The geometric term
\[
-\frac{\dot L}{L} D
\]
produces non-adiabatic transitions proportional to $\dot L/L$.

\medskip

This example demonstrates concretely how a time-dependent domain
is equivalent to a time-dependent self-adjoint extension and
induces a natural measurement structure via quadratic form variation.

\section{Worked Example: Time-Dependent Robin Boundary Condition}

We consider a particle on a fixed interval
\[
\Omega = (0,L), \qquad L>0,
\]
with time-dependent Robin boundary condition at $x=0$ and Dirichlet at $x=L$:

\begin{align}
\partial_x \psi(0,t) + \kappa(t)\psi(0,t) &= 0, \\
\psi(L,t) &= 0,
\end{align}

where $\kappa \in C^1([0,T])$ is real-valued.

\subsection{Self-Adjoint Realization}

Define
\[
H(t)
=
-\frac{\hbar^2}{2m}\frac{d^2}{dx^2}
\]
with domain
\[
D(H(t))
=
\left\{
\psi\in H^2(0,L):
\psi(L)=0,
\;
\partial_x\psi(0)+\kappa(t)\psi(0)=0
\right\}.
\]

\paragraph{Proposition.}
For each $t$, $H(t)$ is self-adjoint and bounded below.

\paragraph{Proof.}
The associated quadratic form is
\[
q_t(\psi)
=
\frac{\hbar^2}{2m}
\int_0^L |\psi'(x)|^2 dx
+
\frac{\hbar^2}{2m}
\kappa(t)|\psi(0)|^2,
\]
with common form domain
\[
Q = \{\psi\in H^1(0,L):\psi(L)=0\}.
\]

The trace map $\psi\mapsto\psi(0)$ is continuous from $H^1(0,L)$ to $\mathbb C$.
Hence $q_t$ is closed and semibounded.
By standard form methods (Kato representation theorem),
$H(t)$ is self-adjoint.
\hfill$\square$

\subsection{Time Derivative of the Quadratic Form}

Since the bulk term is time-independent,
\[
\partial_t q_t(\psi)
=
\frac{\hbar^2}{2m}
\dot\kappa(t)
|\psi(0)|^2.
\]

This formula is exact and requires no boundary flux heuristic.

\subsection{Energy Balance}

Let $\psi(t)$ solve
\[
i\hbar \partial_t \psi = H(t)\psi.
\]

Then
\[
\frac{d}{dt}
\langle \psi(t), H(t)\psi(t) \rangle
=
\partial_t q_t(\psi(t))
=
\frac{\hbar^2}{2m}
\dot\kappa(t)
|\psi(0,t)|^2.
\]

Thus energy variation is entirely localized at the boundary point $x=0$.

\subsection{Induced POVM}

Assume $\dot\kappa(t)\ge 0$ on $\Delta\subset[0,T]$.
Define

\[
M(\Delta)
=
\int_\Delta
U(t,0)^\dagger
\partial_t H(t)
U(t,0)\, dt.
\]

Using quadratic forms:

\[
\langle \psi, M(\Delta)\psi\rangle
=
\frac{\hbar^2}{2m}
\int_\Delta
\dot\kappa(t)
|\psi(0,t)|^2 dt.
\]

Hence:

\begin{itemize}
\item $M(\Delta)\ge 0$ if $\dot\kappa\ge 0$,
\item $M$ is $\sigma$-additive,
\item $M$ defines a POVM measuring boundary amplitude at $x=0$.
\end{itemize}

\subsection{Associated Instrument}

Formally define Kraus density
\[
K(t)
=
\sqrt{\frac{\hbar^2}{2m}\dot\kappa(t)}\,
|\delta_0\rangle\!\langle \delta_0|^{1/2}
\, U(t,0),
\]
where $|\delta_0\rangle$ denotes the boundary evaluation functional
understood via trace embedding.

The instrument is

\[
\mathcal I_\Delta(\rho)
=
\int_\Delta
K(t)\rho K(t)^\dagger dt.
\]

Trace identity:

\[
\mathrm{Tr}[\mathcal I_\Delta(\rho)]
=
\mathrm{Tr}[\rho M(\Delta)].
\]

Thus the modulation of the Robin parameter implements
a continuous measurement of boundary amplitude.

\subsection{Spectral Picture}

Instantaneous eigenvalues $k_n(t)$ solve
\[
k \tan(kL) = -\kappa(t).
\]

Eigenfunctions:
\[
\psi_n(x,t)
=
N_n(t)
\sin(k_n(t)(L-x)).
\]

Time variation of $\kappa(t)$ induces non-adiabatic couplings

\[
\langle \psi_m, \partial_t \psi_n \rangle
=
\dot\kappa(t)
\frac{\partial k_n}{\partial \kappa}
\langle \psi_m, \partial_{k_n}\psi_n \rangle.
\]

Hence transition amplitudes are proportional to $\dot\kappa$.

\subsection{Physical Interpretation}

\begin{itemize}
\item The Robin parameter acts as a boundary coupling constant.
\item $\dot\kappa(t)$ governs energy injection or extraction.
\item The induced POVM measures boundary probability density.
\item The instrument describes continuous boundary monitoring.
\end{itemize}

\medskip

This example demonstrates explicitly that time-dependent self-adjoint
extension parameters correspond to controlled boundary measurement
channels, with statistics determined directly by the quadratic form
variation.

\section{Unified Comparison: Moving Wall vs.\ Robin Modulation}

We now compare the two mechanisms of time-dependent self-adjoint
extensions:

\begin{enumerate}
\item Moving-interval Dirichlet boundary ($L=L(t)$),
\item Fixed-interval time-dependent Robin parameter ($\kappa=\kappa(t)$).
\end{enumerate}

Both are instances of a $C^1$ family of closed quadratic forms
\[
q_t(\psi)
=
\frac{\hbar^2}{2m}\int |\psi'|^2
+
b_t(\psi),
\]
where the bulk kinetic term is time-independent and
all time dependence is localized at the boundary.

\subsection{Quadratic Form Derivatives}

\paragraph{Moving wall.}

After reduction to fixed interval $(0,1)$, one finds
\[
\partial_t q_t(\psi)
=
-2\frac{\dot L}{L^3}
\langle \psi, H_D \psi\rangle
-
\frac{\dot L}{L}
\langle \psi, D \psi\rangle,
\]
where $D$ is the generator of dilations.

Time dependence is global in the bulk but originates from
geometric boundary motion.

\paragraph{Robin modulation.}

For fixed $(0,L)$,
\[
\partial_t q_t(\psi)
=
\frac{\hbar^2}{2m}
\dot\kappa(t)|\psi(0)|^2.
\]

Time dependence is purely local at the boundary point.

\subsection{Structural Comparison}

\begin{center}
\begin{tabular}{l|c|c}
 & Moving Wall & Robin Modulation \\ \hline
Geometry & Domain varies & Domain fixed \\
Reduction & Dilation generator $D$ & No geometric scaling \\
Energy injection & Via $D$ term & Via $|\psi(0)|^2$ \\
POVM density & Boundary gradient & Boundary amplitude
\end{tabular}
\end{center}

Both mechanisms are special cases of:

\medskip
\emph{Time-dependent self-adjoint extension parameter
$\theta(t)$ induces a boundary-local quadratic-form variation
$\partial_t q_t$.}
\medskip

Thus boundary geometry change and boundary coupling modulation
are operationally equivalent from the extension-theoretic viewpoint.

\section{Adiabatic Theorem for Robin Modulation}

We now derive an explicit adiabatic result.

Consider
\[
H(t) = -\frac{\hbar^2}{2m}\frac{d^2}{dx^2}
\]
with boundary condition
\[
\partial_x\psi(0,t)+\kappa(t)\psi(0,t)=0,
\quad \psi(L,t)=0,
\]
where $\kappa \in C^2([0,T])$.

Let $\{\psi_n(t)\}$ be normalized instantaneous eigenfunctions:
\[
H(t)\psi_n(t) = E_n(t)\psi_n(t).
\]

Eigenvalues satisfy
\[
k_n \tan(k_n L) = -\kappa(t),
\quad
E_n(t)=\frac{\hbar^2 k_n(t)^2}{2m}.
\]

\subsection{Spectral Gap Assumption}

Assume a uniform gap condition:

\[
\inf_{t\in[0,T]}
\min_{m\neq n}
|E_m(t)-E_n(t)|
\ge g >0.
\]

\subsection{Derivative of Eigenfunctions}

Differentiate eigenvalue equation:

\[
(H(t)-E_n(t))\partial_t\psi_n
=
(\partial_t E_n - \partial_t H(t))\psi_n.
\]

Taking projection onto $\psi_m$ ($m\neq n$):

\[
\langle \psi_m,\partial_t \psi_n\rangle
=
\frac{
\langle \psi_m,\partial_t H(t)\psi_n\rangle
}{
E_n(t)-E_m(t)
}.
\]

But
\[
\partial_t H(t)
\leftrightarrow
\frac{\hbar^2}{2m}\dot\kappa(t)
|\delta_0\rangle\!\langle \delta_0|.
\]

Thus

\[
\langle \psi_m,\partial_t \psi_n\rangle
=
\frac{\hbar^2}{2m}
\frac{
\dot\kappa(t)
\psi_m(0,t)\psi_n(0,t)
}{
E_n(t)-E_m(t)
}.
\]

\subsection{Adiabatic Transition Amplitude}

Let initial state be $\psi_n(0)$.
Transition amplitude to $m\neq n$ at time $T$ satisfies

\[
|c_m(T)|
\le
\int_0^T
\left|
\langle \psi_m,\partial_t\psi_n\rangle
\right|
dt.
\]

Hence

\[
|c_m(T)|
\le
\frac{\hbar^2}{2m}
\int_0^T
\frac{
|\dot\kappa(t)|
|\psi_m(0,t)\psi_n(0,t)|
}{
|E_n(t)-E_m(t)|
}
dt.
\]

\subsection{Adiabatic Theorem}

\begin{theorem}
Assume:

\begin{enumerate}
\item $\kappa \in C^2([0,T])$,
\item Uniform spectral gap $g>0$,
\item $\sup_t |\dot\kappa(t)| \le \varepsilon$.
\end{enumerate}

Then for initial eigenstate $\psi_n(0)$,
the solution satisfies

\[
\|\psi(t) - e^{i\gamma_n(t)}\psi_n(t)\|
\le
C \frac{\varepsilon}{g},
\]

where $\gamma_n(t)$ includes dynamical and Berry phases,
and $C$ depends on $L$ and boundary trace bounds.
\end{theorem}

\begin{proof}
Use standard adiabatic perturbation expansion
with explicit bound above.
The numerator is $O(\varepsilon)$,
denominator bounded below by $g$.
Integrate over finite time interval.
\end{proof}

\subsection{Berry Phase}

The geometric phase is

\[
\gamma_n^{\text{geom}}(t)
=
i\int_0^t
\langle \psi_n(s),\partial_s\psi_n(s)\rangle ds.
\]

Using previous formula:

\[
\langle \psi_n,\partial_t\psi_n\rangle
=
\frac{\hbar^2}{2m}
\dot\kappa(t)
\sum_{m\neq n}
\frac{
|\psi_m(0)\psi_n(0)|
}{
E_n-E_m
}.
\]

Thus Berry phase is entirely generated by boundary amplitude.

\section{Conclusion of Comparison}

Both moving-wall and Robin-modulation systems exhibit:

\begin{itemize}
\item Self-adjoint extension parameter $\theta(t)$,
\item Quadratic-form derivative localized at boundary,
\item Natural induced POVM,
\item Explicit adiabatic regime controlled by $\dot\theta/g$.
\end{itemize}

The Robin case is analytically cleaner:
domain is fixed and adiabatic estimates are explicit.

The moving-wall case introduces geometric dilation generator
but is structurally equivalent after unitary reduction.

\section{Conclusion}

Time-dependent boundary conditions are rigorously described as
time-dependent self-adjoint extensions.
Their quadratic-form variation generates a natural POVM
encoding boundary energy transfer and defines a quantum instrument.
Thus boundary dynamics admits an intrinsic operational interpretation.

\begin{thebibliography}{99}

\bibitem{Kato}
T. Kato,
\textit{Perturbation Theory for Linear Operators},
Springer, 1966.

\bibitem{Kisynski}
J. Kisy\'nski,
Sur les op\'erateurs de Green des probl\`emes de Cauchy abstraits,
\textit{Studia Math.} 23 (1964), 285--328.

\bibitem{ReedSimon}
M. Reed, B. Simon,
\textit{Methods of Modern Mathematical Physics II},
Academic Press, 1975.

\bibitem{Stinespring}
W. F. Stinespring,
Positive functions on $C^*$-algebras,
\textit{Proc. Amer. Math. Soc.} 6 (1955), 211--216.

\bibitem{vonNeumann}
J. von Neumann,
Allgemeine Eigenwerttheorie Hermitescher Funktionaloperatoren,
\textit{Math. Ann.} 102 (1929), 49--131.

\end{thebibliography}

\end{document}
